\chapter{Glossary}\label{app:glossary}

% [AP-E2/AP-EN] Mirror of anhang/de/C_glossary.tex. Abbreviations + core terms used in Chs. 1--3.
This glossary summarizes the abbreviations and central terms used in the thesis. Each
technical-term entry first names the German source term (the term of the code identifiers and of
the German edition), then the definition and finally the aliases used synonymously in this thesis;
registry and code identifiers are set in typewriter type. The references point to the section that
introduces the term.

\begin{description}
  \item[ABI] Application Binary Interface --- the binary-stable interface through which the builder
  loads the permutation modules (living beings).
  \item[Adaptive vs.\ dynamic vs.\ measurement-driven] Distinction of terms for adapting the
  algorithm behaviour.\ \emph{Adaptive}: adaptation based on run-time-observed data characteristics
  (e.g.\ adaptive node sizes in ART)\@. \emph{Dynamic}: in general usage any decision taken at run
  time; in this thesis, more narrowly, the run-time selection of the fully compiled
  permutation \emph{module} via the binary identity \texttt{binary\_id} (module loading) --- all
  eighteen organ axes are fixed per binary at compile time; in addition there are
  run-time-\emph{variable} parameters of dynamic sub-axes (\emph{Unter-Achsen}: prefetch distance,
  batch size, inline threshold, pool budget; cf.\ Section~\ref{sec:measurement-system}), no
  run-time choice of individual axis strategies.
  \emph{Measurement-driven}: selection of the configuration based on empirically collected
  measurement series --- the leitmotif of this thesis; the corresponding selection mechanism is
  implemented in a first increment (\texttt{best\_binary\_selector}: ranking and delivery of the
  best permutation per metric from the measurement CSVs), the fully automatic,
  load-profile-spanning selection is the subject of the outlook (cf.\ Section~\ref{sec:outlook}).
  \item[Alder Lake] Intel microarchitecture with heterogeneous performance and efficiency cores; the
  second measurement platform of this thesis (prod2, Intel Core i9-12900K) belongs to it. In the
  system registry it is the recombination \texttt{prod2\_alder\_lake} (CPU fabrication
  GenuineIntel/6/151/2, RAM nominal rate 4800~MT/s declared, CAS latency not collected).
  \item[Anatomy] (\emph{Anatomie}) The \emph{body} of a living being: the fixed composition of its
  eighteen axis organs and, above all, their \emph{wiring} --- the generalized interfaces through
  which one organ uses the services of the others (for instance a container family as the node
  substrate of the axis \texttt{node\_type} or as index organization). Not a second model beside
  the search algorithm but its internal structure (cf.\ Section~\ref{sec:anatomy-metaphor}).
  \item[Answer status] (\emph{Antwort-Status}) The level of answering reported in the conclusion
  per main and sub research question: \emph{answered} (a test or a measurement substantiates the
  answer), \emph{structurally answered} (the apparatus that delivers the answer is implemented and
  tested, the measurement is not available) or \emph{open} (neither structure nor measurement is
  available). Without measured values and concrete evidence a question counts as not answered,
  even if its structure is sound (cf.\ Section~\ref{sec:answers}).
  \item[AoSoA] Array of Structures of Arrays --- memory layout that splits records into blocks of
  fixed width and, within a block, lays them out field by field (SoA); it combines the cache-line
  density of the SoA layout with the locality of the AoS layout. In this thesis one value of the
  organ axis \texttt{memory\_layout}.
  \item[Array-indexed page, vector page] (\emph{Punktseite, Kompaktseite}) The two page forms of
  the PRT-ART B$^+$ node by occupancy density: the array-indexed page addresses entries directly
  via an array (8- or 16-bit index, density up to 50~\%), the vector page holds sorted key-value
  vectors (density above 50~\%); the boundary at 50~\% is the page-type boundary, the three
  switching rules of the candidate under test are called \emph{density thresholds}.
  \item[ART] Adaptive Radix Tree~\cite{leis2013art}; reference design with adaptive node sizes
  (Node4/16/48/256).
  \item[Autonomous profile-driven heuristic evaluation]
  (\emph{autonome Heuristik-Profil-Auswertung}) The intended mechanism that automatically derives,
  per load profile, the best axis configuration from the measurement series and stores it as a
  reusable heuristic policy (target mechanism: ML classification of load-profile/axis features onto
  configurations). First building blocks are implemented --- the automatic best-binary selection
  from measurement series (\texttt{best\_binary\_selector}) and the extraction/storage as an XML
  load-profile result ---; the ML classification itself is outlook (cf.\ Section~\ref{sec:outlook}).
  \item[AVX, AVX-512] x86 SIMD extensions (256- and 512-bit, respectively).
  \item[Axis] (\emph{Achse}) An orthogonal partial decision (\enquote{organ}) of a search
  algorithm. The framework assigns its axes to three realms: eighteen organ main axes
  (composition axes in fixed slot order T0--T17, registry label T00--T17) and one optional organ
  meta-meta axis \texttt{disk\_io}; three system main axes (\texttt{target\_isa},
  \texttt{operating\_system}, \texttt{external\_utils}) with their complex wrapper; and the
  measurement realm with the measurement-tooling main axis (cf.\ Section~\ref{sec:axis-triad}).
  Every axis is a main, sub- or meta-meta axis (see the respective entries); only the eighteen
  organ main axes enter the \texttt{binary\_id}.
  \item[Axis success parameter] (\emph{Achsen-Erfolgs-Parameter}) The characteristic quantities
  collected per axis interface, by which an interface call is measured from beginning to end; the
  measuring instrument is the measurement checkpoint \texttt{checkpoint\_measure} (micro level),
  and from the macro level upwards all parameters including the PMC counters are merged
  (cf.\ Section~\ref{sec:measurement-system}). The catalogue of the parameters per axis is not
  available (subject of future work).
  \item[binary\_id] The identity path of a tier binary in the permutation space: the assignment of
  the eighteen organ main axes in canonical order (eighteen segments, the last one
  \texttt{persistence\_target}). System and measurement axes never appear in the
  \texttt{binary\_id} (registry attribute \texttt{binary\_id} with the value \texttt{never}); they
  act via build version, stamp lines and CSV columns.
  \item[built\_stem] Field of the textual plan (revision v1.1) that states, per build step, the
  actually generated file stem --- produced by the very function the build orchestrator calls during
  the build; evidence of the \emph{one} main channel (cf.\ Section~\ref{sec:impl-pipeline}).
  \item[CacheEngineBuilder] (CEB) The autonomous platform-measurement system of the chain and the
  second of the four carrier stages; root of the system realm. The CEB receives from the planner
  the released system-axis space, permutes system and organ axes into tier binaries, builds and
  stamps them per platform and runs the seven-phase pipeline via the \texttt{ExperimentDriver}. Its
  stamp contains the measurement and the system line but no organ line (the static barrier); its
  SHA-512 fingerprint is provenance, not a reuse criterion
  (cf.\ Sections~\ref{sec:mess-chain} and~\ref{sec:impl-pipeline}).
  \item[Candidate under test] (\emph{Prüfling}) The new search algorithm to be integrated on trial
  (here PRT-ART); it is compiled by the builder against the core constituents of the cache engine.
  Attributive short form for constituents of the candidate under test: candidate (candidate slots,
  candidate build axis, candidate variants). This role is distinct from the bare noun: in the
  ranking, the Pareto front, and the selection filter chain, a candidate is a permutation under
  selection (measured in the ranking and the Pareto front, measured or still to be built in the
  filter chain), and the comparison candidates of the workload variation are the compared methods.
  The measurement instrumentation is called \emph{measurement sensor} or, by its literature alias,
  \emph{probe} (see \emph{Measurement sensor}); the fixed names \emph{probe dock} and
  \emph{probe effect} keep their meaning (see \emph{Probe dock}, \emph{Instrumentation artefact}).
  \item[Carrier stage] (\emph{Träger-Stufe}) One of the four stages that carry the anatomy, in
  build order: planner, CacheEngineBuilder (CEB), tier binary, hybrid tier binary (each stage generates
  the next, no hybrid without a tier); in the processing chain the hybrid tier binary sits as the
  optional mux insertion between the CEB probe dock and the tier binary. No stage
  passes through layers --- each contains its layer firmly compiled in; only the configuration
  travels from the user XML to the finished tier binary. Every stage carries a stamp and offers
  three interface surfaces (Figure~\ref{fig:traeger-flaechen}).
  \item[CEB] see \emph{CacheEngineBuilder}.
  \item[checkpoint\_measure] The central measuring instrument of the system: a measurement point
  that collects the input parameters and the time at one instant. On the wall-clock level it is the
  parameter-empty minimum; if the macro or the micro level is active, it additionally performs the
  parameter bracket --- time and observer are united in the checkpoint, the responsibility is
  separated per measurement level. The rows are deposited in the measurement arena, the nesting is
  kept by the stack arena. Alias: measurement checkpoint.
  \item[CLU] Cache-line utilization --- the fraction of used bytes per loaded cache line.
  \item[Complex axis] (\emph{Komplex-Achse}) Bracket over a fixed recombination of several main
  axes that jointly act as \emph{one} axis: the outer bracket over target ISA, operating system and
  extension hardware, and the inner complex axis \texttt{target\_isa} (RAM frequency, CAS latency,
  CPU fabrication); in the identity stamp a complex axis occupies exactly one field
  (cf.\ Section~\ref{sec:complex-axis}). An undeclared member is stamped as undeclared, never
  guessed; the CAS latency is not collected for any measurement machine (n/a: the SMBIOS structure
  type~17 does not carry it).
  \item[Concept] (\emph{Konzept}; code: \texttt{AnatomyGattung}) The topmost level beneath the
  kingdom (level~1): the animal class as the measurement surface --- Map (mammal), Container (bird),
  Graph (plant), HeuristikAdapter (insect) ---, the externally visible measurement interface that
  defines measurability; per family a probe dock materializes it. The concept is not maintained as a
  second datum but determined at run time from the family, via a concept-independent C symbol (the
  symbol bears the same name for every concept). In the product-line reading the concept is the
  feature core composition, the family the feature special composition. Reading note:
  \enquote{Adapter} without prefix always means the container family, the hybrid concept is always
  called HeuristikAdapter. Code bridge: \texttt{AnatomyGattung}, ABI symbol
  \texttt{comdare\_anatomy\_gattung} (renaming in the code after the defence; see
  \emph{Three-level model}, \emph{Family}, \emph{Genus}).
  \item[Conformance gate] (\emph{Konformitäts-Gate}) The unconditional check before every
  measurement: at the probe dock the gate checks the \texttt{std::map} contract of the conformance
  subset (Appendix~\ref{app:interfaces}), so that only comparable things are compared. Under
  \texttt{--debug} it additionally runs the standard test set of the tier interface: the
  family-generic part (valid for all genera) and the genus-specific part (the extended
  special interfaces); this indirect debug mode coincides with the release build and is not kept as a
  separate mode (see \emph{Family}, \emph{Genus}).
  \item[Consumer] (\emph{Konsument}) The role of the thesis code towards the engine: a consumer
  links exclusively against the cache engine and uses one door (\texttt{get\_cache\_engine()}); the
  measurement driver is at the same time its driver. To be distinguished from this is the operator
  --- the human who controls planner and installation. Alias: formal consumer and driver.
  \item[CPI, IPC] Cycles per instruction and instructions per cycle, respectively.
  \item[CRTP] Curiously Recurring Template Pattern --- C++ idiom for static polymorphism: a class
  derives from a class template that receives this very class as its template argument, so that the
  base accesses the derived class at compile time~\cite{coplien1995crtp}. In the framework the
  construction form of the one stamp contract of all carrier stages (the template
  \texttt{StempelBasis} with the derived class and the carrier stage as parameters) and of the axis
  concepts.
  \item[CSV tree] (\emph{CSV-Baum}) The CSV form of the result workbook: not a single child
  document but an ordered tree of CSV files --- one overall folder per workbook, beneath it one
  folder per sheet, therein the CSV files of the individual sub-axis sweeps ---, always generated
  from the xlsx workbook (xlsx to csv, as a strategy of the same factory) and, as a representation,
  its analogously mirrored counterpart; via XML the workbook, the CSV tree or both are deposited at
  the same target. Implementation status: the seam writes the CSV projection as one flat file; the
  tree form is not implemented (Section~\ref{sec:eval-dataflow}).
  \item[Density thresholds] (\emph{Dichte-Schwellen}) The three fixed switching rules of the candidate under test
  PRT-ART between its four page representations by occupancy density, in the order
  Array$_{256}$ (up to 25~\%) $\to$ Array$_{65535}$ (25 to 50~\%) $\to$ Vector$\langle$u8,u8$\rangle$
  (50 to 75~\%) $\to$ Vector$\langle$u16,u16$\rangle$ (above 75~\%); the candidate under test switches by these
  fixed thresholds; the platform-calibrated rules are not available, their determination
  presupposes the measurement runs (sub-question RQ3.a, Section~\ref{sec:rqs};
  Section~\ref{sec:prtart-demo}). The threshold at 50~\% is at the same time the page-type boundary
  between array-indexed and vector page (see \emph{Array-indexed page, vector page}), not a further
  switching rule.
  \item[dTLB] Data translation lookaside buffer; caches data address translations.
  \item[ECDF] Empirical cumulative distribution function --- step function that assigns to each
  value the fraction of observations not exceeding it~\cite{wilk1968probability}; in the
  measurement appendix the representation of the latency distribution across the configurations
  (a single-step staircase for a single configuration).
  \item[Explicitly coded non-observation] (\emph{explizit kodierte Nicht-Erhebung}) Principle of
  the measurement chain: a value not collected is a state of its own (NaN according to IEEE~754 or
  a validity flag~\cite{ieee754_2019}) and is never carried as zero or as a phantom value ---
  neither in the schema nor in CSV, xlsx, table or diagram; a measurement series without a valid
  row produces no output instead of an empty one. The missing value thus remains evaluable as a
  category of its own~\cite{graham2009missing}. Alias: principle of explicitly coded
  non-observation; in the code the vocabulary \texttt{honest-0}/\texttt{honest-empty} of the
  comments and probe-dock messages, in the diagram generator of the tool repository
  (\texttt{diagram\_generator}) the placeholder routine
  \texttt{write\_honest\_empty\_placeholder} (cf.\ Section~\ref{sec:axis-triad}).
  \item[Eytzinger layout] Arrangement of a sorted array in the level order of an implicit complete
  binary tree (children of index~$i$ at $2i$ and $2i+1$), so that binary search proceeds in a
  cache-line- and prefetch-friendly manner~\cite{khuong2017eytzinger}. In this thesis one value of
  the memory-layout axis.
  \item[Family] (\emph{Gattung}; code: \texttt{AnatomyGenus}) The second level of the framework
  (level~2) and the first tier-interface level beneath a concept: the interface that a family carries
  outward as a template, the AND-intersection of the interfaces of all its genera (in the metaphor
  canon the elephant $=$ SearchAlgorithm beneath the mammal Map). Six families: SearchAlgorithm under
  Map, Set, Sequence, Adapter and View under Container (woodpecker, stork, cockatoo, falcon) as well
  as FunctionInterfaceReroute (ant) under HeuristikAdapter. The family is the only maintained source
  of the classification; the concept is determined from it at run time. Per docking family there is
  exactly one probe dock (five docks for SearchAlgorithm, Set, Sequence, Adapter and View), whereas
  the classifying family FunctionInterfaceReroute has none and its binaries dock at the dock of the
  reported target family. The tier binary shows the probe dock its family interface; docks of different
  families are never mixed. The family tests of the conformance gate are generic for all genera of the
  family. Code bridge: \texttt{AnatomyGenus} (renaming in the code after the defence; see
  \emph{Concept}, \emph{Genus}, \emph{Three-level model}).
  \item[Fingerprint] The cryptographic identity hash per carrier stage: the planner carries a
  SHA-256 over version, ISA and operating system (64 hex characters); the tier binary a SHA-512
  over eleven members in fixed order (format~6: format identifier, measurement, system and organ
  line, sub-axis value set, toolchain, enabled-set signature of the build axes, overlay source set,
  measurement gates, hybrid composite, build-version base; 128 hex characters); the CEB a SHA-512
  as provenance. The fingerprint is the store key of stock~1 and the recognition criterion of a
  permutation.
  \item[First-touch policy] NUMA placement rule of the operating system: a memory page is assigned
  to the node whose processor writes it first; therefore the initializing thread decides on the
  memory locality~\cite{lameter2013numa}.
  \item[Fit stamp] (\emph{Passungs-Stempel}) \emph{Design} of a check that validates the
  compiled-in platform \emph{class} of a tier binary against the run-time verdict of the hardware
  detection on every measurement run; not implemented (cf.\ Section~\ref{sec:answers}).
  \item[Forest plot] Representation of several effect estimates with their uncertainty intervals as
  stacked point-and-bar rows~\cite{lewis2001forest}; in the measurement appendix the view of the
  axis exchange per configuration against a reference.
  \item[Full profile, abstract profile] (\emph{Vollprofil, abstraktes Profil}) The two types of
  state-of-the-art profiles: a full profile describes a self-contained search design completely; an
  abstract profile provides only individual organs (OLC as concurrency
  organ~\cite{leis2016olc}, LOUDS as memory-layout organ~\cite{jacobson1989louds}, VAMPIR as
  measurement technology~\cite{berthold2023vampir}) and leaves the remaining ones to the standard
  construction kit --- the same type that the candidate under test PRT-ART corresponds to. Of the 33 profiles,
  30 are full profiles and 3 abstract.
  \item[Genus] (\emph{Genus}) The third level (level~3) and second tier-interface level: the
  subspecies, the implementation of the family interface with possibly extended special interfaces (in
  the metaphor canon the elephant species mammoth $=$ ART, African elephant $=$ HOT, Indian elephant
  $=$ START beneath the family SearchAlgorithm; the giraffe is reserved, CoCo has no code object).
  Beneath the genus lie one or more search strategies (simple $=$ assignment of a single axis,
  complex $=$ composition of several axes); the organs are a classification of their own (level~4)
  that the genus implementation binds. The genus is neither the search-strategy axis nor the
  composition of its organs (the composition from axes is a matter of the version stamp,
  Section~\ref{sec:stamp-model}, not of the taxonomy). The special interfaces of the genus are checked by
  the conformance gate only under \texttt{--debug} (genus tests); to the probe dock the tier binary shows
  its family interface, not the genus. The genus implementation runs, per function, one main algorithm
  that works exclusively through axis-interface calls, and it calls foreign axes first. As-is: today the code carries the genus only as an axis assignment; the genus
  implementation class is missing (an acknowledged architecture error, fix after the defence). The
  hybrid artefact is classified by the family FunctionInterfaceReroute; \texttt{genus()} of a hybrid
  binary, by contrast, reports the inherited target family, the rerouting remains transparent to the
  outside. Code bridge: \texttt{AnatomyGenus} denotes in the code the \emph{family} of this thesis,
  not the genus (renaming after the defence; see \emph{Family}).
 
  \item[Gitlink] The commit of a submodule recorded in the superordinate repository (entry mode
  160000); \texttt{.gitmodules} declares per submodule the path, the URL and optionally the
  branch~\cite{git2026gitmodules}. The super-repository thus carries the authoritative state of
  engine, candidate under test and manuscript.
  \item[GoF] \enquote{Gang of Four} --- the canonical design patterns after Gamma et
  al.~\cite{gamma1994patterns}.
  \item[golden reference catalogue] (\emph{golden-Referenzkatalog}) The fixed regression detector
  of the whole system: seventeen organ axes with two values each, \texttt{persistence\_target}
  fixed to the in-memory representative, yield $2^{17} = 131072$ identities, anchored via a CRC-64
  checksum instead of via the materialized list --- under the simplifying, superseded assumption of exactly
  two configurations per axis; the axes carry several configurations, the catalogue remains the regression
  detector, not an upper bound. It is to be kept apart from the upper bound of binary identities per system
  permutation: the product of the configurations selected per axis in the XML experiment, which the planner
  computes dynamically per profile and platform from the activated catalogues and which is known only at
  execution time (cf.\ Section~\ref{sec:mess-chain}).
  \item[H2 dossier] (\emph{H2-Akte}) The tool-computed data source of the evaluation hypothesis H2
  (code quality per source): a weighted cppcheck finding density per kLOC over the vendored
  original sources of the state-of-the-art designs; without a vendored original source the dossier
  carries \texttt{n/a} with a reason. Not to be confused with the PRT-ART telemetry metrics
  H1--H3.
  \item[HBM] High-Bandwidth Memory.
  \item[HDR histogram] High dynamic range histogram with logarithmically graded bins and
  guaranteed relative precision across many orders of magnitude~\cite{tene_hdrhistogram}; in the
  evaluation the distribution view of the latencies, whereas the percentiles are determined
  nearest-rank from the raw values (see \emph{nearest-rank median}).
  \item[header-only] Library form whose entire code lies in header files and is compiled upon
  inclusion; no separate library artefact arises. The candidate under test PRT-ART and the hybrid stage exist in
  this form.
  \item[HOT] Height Optimized Trie~\cite{binna2018hot}.
  \item[HybridAware] The sixth core mode of the measurement model and value \texttt{HybridAware}
  of the scheduling dimension \texttt{hetero\_core\_dispatch} (beside \texttt{None},
  \texttt{PCoresOnly} and \texttt{ECoresOnly}): the run knows the heterogeneous core topology, pins
  its threads at run time via the affinity interface of the operating system onto selected cores of
  both types and measures there with the counter set belonging to the core type; precondition of
  the cache-line-awareness measurement and therefore a mandatory part of the measurement model
  (cf.\ Section~\ref{sec:axis-triad}).
  \item[Hybrid tier binary] The fourth carrier stage (level~5b: inserted between the CEB probe dock and
  the tier binaries as a multiplexer with slots) in the build order, in the processing chain the
  optional mux insertion between the CEB probe dock and the tier binary: a module of the concept
  HeuristikAdapter
  (family FunctionInterfaceReroute) that inherits at compile time the interfaces of one concept and
  one family and puts calls through to attached tier binaries according to heuristic criteria; it
  exports the same mandatory symbols as a tier, reports the target family to the outside and is
  built exclusively by means of the CEB\@. Implementation status: classification, dock layer, XML
  parser and binary proxy are implemented; the module export as a separate compiled artefact, the
  reroute router, the displacement, the state synchronization (memento sidecar) and the live
  ranking are not implemented (cf.\ Section~\ref{sec:impl-hybrid-lager}).
  \item[Instrumentation artefact] (\emph{Instrumentierungs-Artefakt}; literature term: probe
  effect) The alteration of the observed behaviour by the measuring instrument
  itself~\cite{gait1986probe}; in this thesis the difference between the tier binary with and
  without measurement sensors, which is checked via the plausibility invariant (slower with
  measurement sensors than without) and factored out via the measurement-error interpolation.
  \item[Instrumentation shut-off level] (\emph{Abschaltungsstufe}) Graduation in which the
  measurement instrumentation of a tier binary is withdrawn: level~1 switches off the collection
  of the axis parameters in the macro and the micro level, so that the wall-clock measurement
  remains for \texttt{measure} and \texttt{compare}; level~2 (permutative, \texttt{release}) also
  removes the wall-clock measurement and with it the measurement surfaces and interfaces from the
  compiled artefact. Implementation status: the code contains the measurement gates G2 and G3,
  whose state is a fingerprint member, and the release mode without instrumentation; the two-level
  shut-off as a permutation quantity is intended but not implemented
  (Figure~\ref{fig:abschaltungsstufen}, Section~\ref{sec:impl-system-axes}).
  \item[Interface surface] (\emph{Interface-Fläche}) One of the three contract surfaces that every
  carrier stage offers: surface~1 the family interface (run time, Abstract Factory), surface~2 the
  stamp (compile time, ABI-stable), surface~3 the measurement cut-through (\texttt{IMessVisitor}),
  which crosses the measurement as a seam of its own so that concept and family interfaces remain
  unchanged (Figure~\ref{fig:traeger-flaechen}). Alias: contract surface, carrier surface.
  \item[IQR] Interquartile range --- difference between the third and the first quartile (p75 $-$
  p25); in the measurement appendix the dispersion measure of the sibling-pair differences in the
  exchangeability long tables and the forest plots (whisker = IQR); a box plot of the latencies is
  deliberately not drawn there, because the schema carries per operation only the percentiles p50,
  p99 and p999 but no quartile boundaries (p25/p75)~\cite{mcgill1978boxplots}.
  \item[Living being] (\emph{Lebewesen}) A completely composed tier module of a concept --- all
  families and genera are living beings; the family SearchAlgorithm is the case with the full eighteen-axis anatomy,
  and only the ABI adapter view narrows the term to it (technical identifier there: SearchAlgorithm
  instance). The user XML carries a living being with one or more tiers. Delimitation: axis-less
  Graph modules are viruses (see \emph{Virus}).
  \item[LOUDS] Level-Ordered Unary Degree Sequence --- succinct (space-saving close to the
  information-theoretic bound) tree encoding that stores the node degrees level by level in
  unary~\cite{jacobson1989louds}.
  \item[M0--M3] (meta-levels M0--M3) The four abstraction levels by which this thesis orders its
  axis terms, following the four-level meta-level architecture of the Object Management Group
  (MOF)~\cite{omg2016mof}: M3 the meta-meta axes as the indirection of a description; M2 the meta
  axes as the literal classification of the main and sub-axes carried under their term; M1
  conception and run-time implementation; M0 the binary code. The assignment of M3 is fixed, the
  cut between M2 and M1 is a reading (Figure~\ref{fig:m0-m3}, Section~\ref{sec:impl-system-axes}).
  \item[Main axis] (\emph{Haupt-Achse}) A compile-time-static axis whose choice is compiled firmly
  into the CEB or the tier binary and appears in the build legend: the eighteen organ main axes
  (identity-forming), the three system main axes and the measurement-tooling main axis. Its
  sub-axes are permuted by the planner at run time. In the stamp entry the main axis is the
  bracket-free level~0.
  \item[Measurement level] (\emph{Mess-Ebene}) One of the three levels of the measurement-tooling
  main axis, in exactly this order: \emph{wallclock} (w) --- the wall clock of the overall runs in
  the CacheEngineBuilder across the load frameworks; \emph{macro} (ma) --- the time of the
  family-interface functions of the tier binary; \emph{micro} (mi) --- time and success parameters
  per axis interface. Subsets are permitted, any other order is syntactically wrong
  (cf.\ Section~\ref{sec:measurement-system}). Alias: w/ma/mi; wallclock, macro, micro level.
  \item[Measurement realm] (\emph{Mess-Realm}) The third axis realm beside organ and system axes:
  the concerns of measuring itself, rooted at the experiment \emph{planner}. Its main axis is the
  measurement tooling (wallclock/macro/micro), the fan-out of the CEB type; its sub-axes are the
  sixteen measurement categories (CSV columns), the work mode, workloads and datasets as well as
  the write-back methods; the load framework is its meta-meta main axis. All measurement-realm
  axes are strictly binary-id-neutral (cf.\ Section~\ref{sec:axis-triad}).
  \item[Measurement sensor] (\emph{Messfühler}) The measurement instrumentation compiled into a
  tier binary (measurement cut-through, observers, PMC counters); its assignment is encoded in the
  measurement-gates member of the fingerprint, so that the binary with and the binary without
  measurement sensors possess different identities and both lie in the store. The difference of
  their run times is the instrumentation artefact. The word \emph{probe} is used in this thesis as
  an alias of the measurement sensor; the candidate under test is never called \emph{probe} (see
  \emph{Candidate under test}).
  \item[Merge strategy] (\emph{Verbund-Strategie}) The manner, declared per axis, in which a candidate
  slot is joined with the engine axis: \texttt{Verbund1\_CeOnly} (only the engine, no candidate
  directive), \texttt{Verbund2\_Replace} (the candidate axis replaces the engine axis, with fallback;
  the default), \texttt{Verbund2\_Hybrid} (engine and candidate under test hybrid per candidate
  under test) and
  \texttt{Verbund3\_Union} (the non-redundant union); the XML tokens are \texttt{replace},
  \texttt{merge} and \texttt{union}. Alias: stage~1 to~3 (old name), full join (for
  \texttt{union}).
  \item[MESI, MOESI, MESIF] Cache coherence protocols (base protocol and AMD/Intel variants,
  respectively).
  \item[Meta-meta axis] (\emph{Meta-Meta-Achse}) An axis of level M3: it describes not an
  assignment but the indirection of a description, and stamps as the bracket appendage of its realm
  line; the meta-meta level of a stamp entry is its bracket depth. Each realm has its own
  discriminator: \texttt{organ\_meta\_meta} (\texttt{disk\_io}), \texttt{system\_meta\_meta} (the
  hardware-family axes beneath the hub of the extension hardware, such as the SIMD family) and
  \texttt{measurement\_meta\_meta} (\texttt{load\_framework}); the PMC collection counts among the
  meta-meta measurement axes.
  \item[Metaphor canon] (\emph{Metapher-Kanon}) The biological reading of the architecture, used
  uniformly in this thesis: kingdom $=$ Animalia (\texttt{IAnatomyBase}); animal class $=$ concept
  (\enquote{mammal} $=$ Map, \enquote{bird} $=$ Container, \enquote{plant} $=$ Graph --- botanically
  the kingdom Plantae, the metaphor deliberately reaches beyond the kingdom Animalia here ---,
  \enquote{insect} $=$ HeuristikAdapter); species $=$ family (\enquote{elephant} $=$ SearchAlgorithm,
  \enquote{woodpecker} $=$ Set, \enquote{stork} $=$ Sequence, \enquote{cockatoo} $=$ Adapter,
  \enquote{falcon} $=$ View, \enquote{ant} $=$ FunctionInterfaceReroute); subspecies $=$ genus
  (\enquote{mammoth} $=$ ART, \enquote{African elephant} $=$ HOT, \enquote{Indian elephant} $=$
  START); organ $=$ axis; anatomy $=$ the body of a living being; animal (\emph{Tier}) $=$ a living
  being of a family, the individual its tier binary; virus $=$ axis-less Graph module; blood $=$
  the liquid organ that connects all organ axes and is condensed in the planner into one
  compile-time configuration; habitat $=$ the system axes; laboratory $=$ the measurement axes
  (cf.\ Section~\ref{sec:anatomy-metaphor}). Where the text means the
  measurement cross-section of the host, it says so without metaphor.
  \item[MLP] Memory-level parallelism --- number of simultaneously outstanding memory accesses that
  a core can overlap; bounded by the miss status holding registers (see \emph{MSHR}).
  \item[MPKI] Misses per kilo-instructions --- cache or TLB misses per thousand executed
  instructions; normalizes miss counts across programs of different length.
  \item[MSHR] Miss Status Holding Register --- buffer of a cache that manages outstanding misses
  and thereby bounds the maximum number of simultaneously open accesses (MLP).
  \item[nearest-rank median] Percentile definition without interpolation: the $p$-quantile is the
  value at rank $\lceil p \cdot n \rceil$ of the sorted sample (type~1 according to Hyndman and
  Fan~\cite{hyndman1996quantiles}); all percentiles of the evaluation are determined this way from
  the raw values, the HDR distribution complements the view.
  \item[NINE] Non-Inclusive-Non-Exclusive --- inclusion policy of a shared last-level cache that
  guarantees neither inclusion nor exclusion with respect to the private caches.
  \item[NUMA] Non-Uniform Memory Access.
  \item[OLC] Optimistic Lock Coupling --- optimistic concurrency technique (ARTSync,
  P08)~\cite{leis2016olc}.
  \item[Organ] Metaphorical term for an \emph{axis} (level~4 beneath the three-level model); a search
  algorithm possesses eighteen organ axes. Beside the eighteen organ main axes the organ realm
  carries an optional organ meta-meta axis (\texttt{disk\_io}), which takes effect only with a disk
  write-back target (see \emph{Axis}, \emph{Meta-meta axis}).
  \item[Overlay source set] (\emph{Overlay-Quellmenge}) The one set of source files over which the
  overlay member of the tier fingerprint hashes: per axis the files of its own implementation plus
  the shared tier substance, concatenated in fixed order per axis category and hashed once with
  SHA-512; the same set is the population of the version lock.
  \item[Page type] (\emph{Seitentyp}) The candidate slot \texttt{page\_type} (prt-art
  \texttt{axis\_01}): the node form of the PRT-ART B$^+$ tree; in the canon of the engine it
  corresponds to the organ main axis \texttt{node\_type} (slot T04). The slot numbers of the candidate under test
  do not follow the slot order T00--T17 of the engine.
  \item[Pareto front] Set of the configurations that cannot be improved in any metric without
  deteriorating in another (Pareto dominance)~\cite{emmerich2018multiobjective}; in the measurement
  appendix the view of the latency-throughput-memory trade-off of the permutations.
  \item[P-core, E-core] (\emph{performance core, efficiency core}) The two core types of hybrid
  processors (on Intel since Alder Lake); the measurement pins threads per core type and collects
  the PMC counters separately. The core dispatch is the dimension \texttt{hetero\_core\_dispatch}
  of the scheduling sub-axis at \texttt{target\_isa}, the core class the run-time sub-axis
  \texttt{core\_class}; the measurement model carries six core modes (four pinning modes per core
  type, the unpinned reference run and HybridAware as the sixth mode), the connection to the run
  profile is not implemented (cf.\ Section~\ref{sec:axis-triad}).
  \item[Pitchfork layout] Convention for the directory structure of native C/C++ projects
  (\texttt{apps/}, \texttt{libs/}, \texttt{tools/}, \texttt{tests/}, \texttt{docs/},
  \texttt{external/} --- in the engine \texttt{ext/}), which the engine
  follows~\cite{pike2018pitchfork}.
  \item[Planner] (\emph{Planer}) The experiment planner (\texttt{ExperimentPlanDirector}):
  resolves the user XML against the three offering catalogues, selects the measurement axes and steers
  the builder orders of the chain; root of the measurement realm and first of the four carrier
  stages, as a tool a subcommand program (\texttt{validate}, \texttt{plan}, \texttt{cache-key},
  \texttt{fingerprint}, \texttt{status}, \texttt{check-size}, \texttt{simulate}, \texttt{version},
  \texttt{help}); its stamp is the only one to carry the version line and a SHA-256 fingerprint
  (cf.\ Section~\ref{sec:mess-chain}).
  \item[PMC] (Hardware) Performance Monitoring Counter. In the measurement realm the counter path
  is a building block of the collector axis (regime PmcCounter); the model-bound raw events are
  supplied by a catalogue per microarchitecture, and in the axis typology PMC counts among the
  meta-meta measurement axes.
  \item[Probe dock] (\emph{Prüf-Dock}) The builder-side measurement orchestration per family: the
  concept defines the measurability, per family a probe dock materializes it. Exactly one dock per family (SearchAlgorithm, Set, Sequence, Adapter, View) loads a living being,
  drives it through the interface operations and measures it; hybrid docks bind attached tiers.
  The dock is no ABI boundary --- the ABI boundary remains the family-specific drive sub-interface
  including the POD snapshot; the persistence of the results is taken over by the result intake
  into the experiment tree and the write-back (cf.\ Section~\ref{sec:anatomy-levels}).
  \item[Provenance level] (\emph{Provenienz-Stufe}) Trust level of a collected hardware value ---
  configured-measured, SPD-JEDEC base value or declared-not-measured --- carried by every result
  together with its collection timestamp; a lower level displaces a higher one only with the label
  travelling along (cf.\ Section~\ref{sec:hw-probe}).
  \item[PRT-ART] Probst Redirect Tree / ART hybrid --- the candidate under test of this thesis.
  \item[Rank] (\emph{Rang}) Centrality classification of the state-of-the-art publications
  (rank-1/2/3); explained in Section~\ref{sec:rqs}. Not to be confused with the ranking of the
  measurement (see \emph{Ranking}).
  \item[Ranking] (\emph{Rangbildung}) The ordering of the measured permutations by a metric across
  the three benchmarking granularities (thesis assignment, sub-task~7); the first increment is the
  ranking per metric from the measurement CSVs with delivery of the best binary
  (\texttt{best\_binary\_selector}).
  \item[RCU] Read-Copy-Update --- reclamation scheme (McKenney, P29)~\cite{mckenney2001rcu}.
  \item[Realm] One of the three domains of the axes, each with its own registry, registry XML and
  stamp line: the organ realm (root: the anatomy of the tier), the system realm (root: the CEB)
  and the measurement realm (root: the planner); the canonical order is measurement, system,
  organ axes. Alias: axis home, axis realm.
  \item[Registry XML] The machine-readable image of the axis inventory (axes, building blocks,
  sub-axes, machine signatures) generated per realm by compile-time reflection; it is never edited
  by hand and is the offer against which the planner resolves the user XML\@.
  \item[REUSE, SPDX] REUSE specification (Free Software Foundation Europe): convention of
  assigning every file of a repository to a licence via SPDX licence identifier and
  \texttt{REUSE.toml}; SPDX is the standardized identifier catalogue for licences. The assignment
  re-licenses nothing --- the licence files remain authoritative~\cite{fsfe2024reuse}.
  \item[RRIP, DRRIP] (Dynamic) Re-Reference Interval Prediction --- cache replacement
  heuristic~\cite{jaleel2010rrip}.
  \item[SIMD] Single Instruction, Multiple Data.
  \item[SOSD] Search On Sorted Data --- index benchmark~\cite{kipf2019sosd}.
  \item[SOTA] State of the Art.
  \item[SPD, JEDEC] Serial Presence Detect --- the identification memory of a memory module
  standardized by JEDEC (the standards body of the memory industry), which carries, among other
  things, the JEDEC base frequency~\cite{jedec_jesd400_5}; in the hardware detection the middle
  provenance level (SPD-JEDEC base value) between configured-measured and declared-not-measured
  (see \emph{Provenance level}).
  \item[Stamp] (\emph{Stempel}) The compiled-in, ABI-stable identity inscription of a carrier: per
  realm one line of entries of the form \texttt{achse=algo@X.Y.Z} with optional flags and
  recursive bracket groups (\texttt{c\{p.e\}} is \texttt{c} with the sub-members \texttt{p} and
  \texttt{e}, never a flat abbreviation), plus version line and fingerprint. Four carrier stages
  and seven stamp interfaces form a three-valued admission matrix of 28 cells (mandatory,
  permitted, forbidden); the CEB carries no organ line, only the planner a version line, only the
  hybrid attached tiers. The stamp is at once identity, store key and recognition criterion; the
  version expression of the driver, by contrast, serves only classification
  (Figure~\ref{fig:traeger-flaechen}). Alias: identity stamp, stamp system.
  \item[Stock] (\emph{Bestand}) One of the four deposit classes of the store, fixed
  compile-time-statically in the key schema: stock~1 the built tier binaries, keyed by their
  SHA-512 fingerprint, the measurement cell bracketed beside it; stock~2 the measured values, keyed
  by the fingerprint of the measured binary and the hardware identity of the measuring machine, the
  platform cell bracketed beside it and never fused into the digest; stock~3 the synthesis batches
  with the synthesized measurement curves as files; stock~4 the solutions per user XML
  (cf.\ Section~\ref{sec:lager}).
  \item[Store] (\emph{Lager}) The persistent holdings of the apparatus, in which what is built and
  what is measured lies under stamp keys (see \emph{Stock}): the store tree is organized by
  concept, family and realm, a permutation is recognized via its fingerprint instead of being built
  anew, measurement is always switched on and re-measures in the store only what is missing there.
  Every machine configuration has its own measurement stock; an algorithm version change triggers
  the rebuild of the affected store stack together with re-measurement. Implementation status:
  stocks~1 and~2 as well as the key schema of all four stocks are implemented; the mechanics of
  stocks~3 and~4, the curve facet, the storage path of the hybrid tier and the target selection
  via the XML form the full store build-out and are not implemented
  (Figure~\ref{fig:lager-bestaende}, Section~\ref{sec:lager}).
  \item[Sub-axis] (\emph{Unter-Achse}) An axis beneath a main axis, as a rule dynamic and permuted
  by the planner at run time: its values travel as settings via the probe dock into the tier
  binary, appear in no build legend and never in the \texttt{binary\_id}. Examples: the run-time
  parameters prefetch distance, batch size, inline threshold, pool budget and thread count as well
  as the sixteen measurement categories, the workload and the work mode (measurement realm);
  optimization level, atomics capability, SIMD family, NUMA node, page size and core class (system
  realm). Individual sub-axes are compile-time-static tuples of their main axis (for instance
  scheduling and compiler) and act via the build version. Registry token \texttt{sub\_axis}.
  Alias: subaxis. Sub-axis and permutation axis are not mutually exclusive.
  \item[Subset] (\emph{Teilmenge}) In this thesis in two fixed meanings: the permissible subset of
  the measurement levels (any subset of wallclock/macro/micro in canonical order) and the
  conformance subset of the interface operations, which the conformance gate checks before every
  measurement (under \texttt{--debug} additionally the standard test set of the tier interface; Appendix~\ref{app:interfaces}).
  \item[Subsystem] In this thesis the four carrier subsystems of the chain (measurement driver,
  CacheEngineBuilder, CacheEngine, candidate under test); in the code, by contrast, the twelve
  \texttt{ISubEngine} stages of the recommendation pipeline behind the facade
  \texttt{ICacheEngine}. The two meanings are not mixed in the text.
  \item[SuRF] Succinct Range Filter~\cite{zhang2018surf}.
  \item[SVE, NEON, RVV] ARM (SVE, NEON) and RISC-V (RVV) SIMD extensions, respectively.
  \item[Three-level model] (\emph{Drei-Ebenen-Modell}) The three interface levels of the one
  architecture: \emph{concept} (level~1, the animal class as the measurement surface --- Map,
  Container, Graph, HeuristikAdapter), \emph{family} (level~2, the tier interface as a template ---
  SearchAlgorithm, Set, Sequence, Adapter, View, FunctionInterfaceReroute) and \emph{genus} (level~3,
  the subspecies as the implementation of the family interface --- ART, HOT, START). Beneath them lie
  the \emph{organs}/\emph{axes} (level~4) as a classification of their own, above them the kingdom
  (level~0); tier binary (level~5a) and hybrid tier (level~5b) are the post-phase of the compilation.
  See Section~\ref{sec:anatomy-levels}.
  \item[Tier binary] The third carrier stage (level~5a of the concept hierarchy: the post-phase of the
  compilation, the individual): the dynamically loadable module of a tier, built by
  the CEB per permutation, which contains exactly one organ assignment (\texttt{binary\_id})
  together with the system and measurement choice firmly compiled in, is loaded via
  letter-identical C symbols (the four primal symbols, the concept and family query and the
  version lines) and lies in the store as stock~1. Alias: tier, permutation module. The word \emph{Tier} is German for
  \emph{animal} and names, in the metaphor canon of this thesis, the fully composed module --- the living
  being with all its organs ---, not a level or layer.
  \item[TLB] Translation Lookaside Buffer.
  \item[Two-gate protocol] (\emph{Zwei-Gate-Protokoll}) The acceptance protocol per implementation
  package: gate~1 is the double-run gate of the engine repository --- the full test suite runs in
  four cells (gcc and clang, each as release and debug build) without change skip, and only a run
  without failed tests is accepted, without an exception list; gate~2 is the integration gate of
  the project repository (super-repository), whose pipeline builds error-free without an exception
  set (cf.\ Section~\ref{sec:impl-contracts}).
  \item[two\_phase\_valid] Validity flag of a measurement row: only a measurement whose workload
  was run through the two-phase load profile (build-up, then measurement phase) counts as valid;
  rows with \texttt{two\_phase\_valid} equal to \texttt{false} do not enter median, matrix or
  diagram.
  \item[Version lock] (\emph{Versions-Lock}) The checked-in lock file that records category, axis
  version and SHA-256 per source file of the overlay source set; every change to a recorded file
  demands the deliberate regeneration commit, a deviation against the discovery is a failure, and
  a CI gate checks it.
  \item[Virus] An axis-less Graph or pure-mathematics module at the same root as the living beings:
  without axis system, without composition and outside the permutation tree, but measurable with
  its own flat measurement POD (latency, throughput, module-specific counters). Plural: viruses.
  \item[working\_set\_n] The column of the working-set size (number of keys) of a measurement
  series; it is the abscissa of the sweep curves of the curve loader, and the cell tokens
  \texttt{n/a} and \texttt{failed} are not numbers --- without a measured value no support point
  arises.
  \item[Workload selection bias] (\emph{Lastprofil-Selektionsverzerrung}) The distortion of a
  comparison by choosing the load profiles such that one design performs preferentially well under
  them~\cite{blackburn2016truth}; the thesis counters it with the global standard workload as a
  permutation over all available workloads and with profiles reproduced per publication.
  \item[Work mode] (\emph{Arbeitsmodus}) The four modes of the run chain, in this order:
  \texttt{build} (building the tier binaries), \texttt{measure} (the deterministic measurement
  run), \texttt{compare} (comparison of the measurement levels and permutations from the
  measurement store) and \texttt{release} (delivery without measurement sensors); the mode-specific procedure of
  \texttt{compare} and \texttt{release} is envisioned for later main waves. The work mode is
  a sub-axis of the measurement realm (\texttt{work\_mode}), not a property of the binary
  identity; the mode token \texttt{compare} is to be distinguished from the identically named
  alias of the wallclock level in the measurement configuration. Alias: mode, run chain.
  \item[XML load-profile result] (\emph{XML-Lastprofil-Ergebnis}) The persistence format demanded
  by the thesis assignment (sub-task~8, automation of system detection and binary selection) in
  which the heuristic policy extracted for an overall algorithm --- reusable per architecture focus
  --- is stored under reference to the applied load profiles; it serves as input to the autonomous
  profile-driven heuristic evaluation. Import parser and export writer including the extractor are
  implemented (round-trip-capable schema \texttt{comdare\_load\_profile}); outlook is the
  autonomous heuristic evaluation built on top (ML classification).
  \item[YCSB] Yahoo!\ Cloud Serving Benchmark~\cite{cooper2010ycsb}.
  \item[Zen~5] AMD microarchitecture (CPUID family 26); the first measurement platform of this
  thesis (prod1, AMD Ryzen~9~9950X3D, Granite Ridge, two core complexes) belongs to it, and its % chktex 29
  raw-event catalogue is cross-checked by measurement. In the system registry it is the
  recombination \texttt{prod1\_zen5}. In this thesis the term denotes exclusively the
  microarchitecture.
\end{description}
